% RQ5 Methodology: Inter-DAO Cooperation Experiments

\subsection{Experimental Design}

We simulate five cooperation scenarios varying DAO characteristics and cooperation structures:

\begin{table}[htbp]
\centering
\caption{RQ5 Experiment configuration (Experiment 07)}
\label{tab:rq5-experiment}
\begin{tabular}{@{}llll@{}}
\toprule
Scenario & DAOs & Structure & Key Variable \\
\midrule
Homogeneous & 3 $\times$ 200 members & Symmetric & Baseline cooperation \\
Heterogeneous & 500, 200, 100 & Diverse sizes & Complementarity effects \\
Asymmetric & 1000, 100 & Large-small pair & Fairness mechanisms \\
Hub-spoke & 1 hub + 4 spokes & Central coordinator & Coordination efficiency \\
Competitive & 3 $\times$ 200 & Rival DAOs & Cooperation despite competition \\
\bottomrule
\end{tabular}
\end{table}

Each scenario runs 60 times with different seeds.

Total runs: $5 \times 60 = 300$

\subsection{Scenario Details}

\subsubsection{Scenario 1: Homogeneous DAOs}

Three DAOs with similar characteristics:
\begin{itemize}
    \item 200 members each
    \item Similar preference distributions
    \item Equal bargaining power
    \item 20\% membership overlap between pairs
\end{itemize}

Hypothesis: High cooperation success but modest surplus.

\subsubsection{Scenario 2: Heterogeneous DAOs}

Three DAOs with diverse profiles:
\begin{itemize}
    \item Large DAO (500 members), Medium (200), Small (100)
    \item Different specializations (technical, community, treasury)
    \item Complementary capabilities
    \item 10\% membership overlap
\end{itemize}

Hypothesis: Higher potential surplus but coordination challenges.

\subsubsection{Scenario 3: Asymmetric Pair}

One large DAO cooperating with one small DAO:
\begin{itemize}
    \item Large: 1,000 members
    \item Small: 100 members
    \item Power imbalance
    \item Test fairness mechanisms (equal split vs. Shapley)
\end{itemize}

Hypothesis: Explicit fairness mechanisms required for stability.

\subsubsection{Scenario 4: Hub-Spoke}

One central ``hub'' DAO coordinating with four ``spoke'' DAOs:
\begin{itemize}
    \item Hub: 500 members, coordination role
    \item Spokes: 100 members each, specialized
    \item Hub facilitates inter-spoke cooperation
\end{itemize}

Hypothesis: Hub improves coordination efficiency.

\subsubsection{Scenario 5: Competitive Cooperation}

Three DAOs that compete but may cooperate on public goods:
\begin{itemize}
    \item 200 members each
    \item Competing for users/market share
    \item Shared interest in ecosystem infrastructure
    \item Mixed incentives
\end{itemize}

Hypothesis: Cooperation limited to clear public goods.

\subsection{Cooperation Mechanisms Tested}

\begin{itemize}
    \item \textbf{Joint proposals}: Multi-DAO approval requirements
    \item \textbf{Resource pooling}: Shared treasury contributions
    \item \textbf{Shapley fairness}: Division proportional to marginal contribution
    \item \textbf{Equal split}: Simple equal division regardless of size
\end{itemize}

\subsection{Hypotheses}

\begin{itemize}
    \item \textbf{H5.1}: Homogeneous DAOs achieve higher cooperation success rates
    \item \textbf{H5.2}: Heterogeneous DAOs generate larger cooperation surplus
    \item \textbf{H5.3}: Asymmetric relationships require explicit fairness mechanisms
    \item \textbf{H5.4}: Hub-spoke structures improve multi-DAO coordination efficiency
    \item \textbf{H5.5}: Overlapping membership increases cooperation probability
\end{itemize}

\subsection{Metrics}

Primary metrics for RQ5:

\begin{description}
    \item[Cooperation Success Rate] Joint proposals that pass all participating DAOs
    \item[Surplus Generated] Value created by cooperation minus standalone values
    \item[Resource Flow Volume] Treasury transfers between DAOs
    \item[Alignment Score] Correlation of voting patterns across DAOs
    \item[Coalition Stability] Duration of cooperation relationships
\end{description}

\subsection{Statistical Analysis}

\subsubsection{Scenario Comparison}

ANOVA comparing cooperation success across scenarios.

\subsubsection{Mechanism Effectiveness}

Compare fairness mechanisms (Shapley vs. equal split) in asymmetric scenarios.

\subsubsection{Correlation Analysis}

Relationship between membership overlap and cooperation success.
